\section{Lecture 2: Matrix Formulation and Multiple Regression}

\subsection{Matrix Form of Simple Linear Regression}

The simple linear regression model is
\[
Y_i = \beta_0 + \beta_1 x_i + \varepsilon_i,
\qquad \varepsilon_i \sim \mathcal{N}(0,\sigma^2),
\quad i = 1,\dots,n.
\]

In vector form,
\[
\mathbf{Y}
=
\begin{pmatrix}
Y_1\\
\vdots\\
Y_n
\end{pmatrix}
=
\begin{pmatrix}
1 & x_1\\
\vdots & \vdots\\
1 & x_n
\end{pmatrix}
\begin{pmatrix}
\beta_0\\
\beta_1
\end{pmatrix}
+
\begin{pmatrix}
\varepsilon_1\\
\vdots\\
\varepsilon_n
\end{pmatrix}.
\]

Define
\[
\mathbf{X}
=
\begin{pmatrix}
1 & x_1\\
\vdots & \vdots\\
1 & x_n
\end{pmatrix},
\quad
\bbeta
=
\begin{pmatrix}
\beta_0\\
\beta_1
\end{pmatrix},
\quad
\beps
=
\begin{pmatrix}
\varepsilon_1\\
\vdots\\
\varepsilon_n
\end{pmatrix}.
\]

The model can be written as
\[
\mathbf{Y} = \mathbf{X}\bbeta + \beps,
\qquad
\beps \sim \mathcal{N}(\mathbf{0}_n, \sigma^2 \mathbf{I}_n).
\]

\subsection{Multiple Linear Regression Model}

\textbf{Model:}
\[
Y_i = \beta_0 + \beta_1 x_{i1} + \cdots + \beta_p x_{ip} + \varepsilon_i,
\qquad i = 1, 2, \ldots, n
\]

\textbf{Parameters:}
\[
\beta_0, \beta_1, \ldots, \beta_p
\]

\textbf{Estimates of the parameters:}
\[
\bhat{\beta}_0, \bhat{\beta}_1, \ldots, \bhat{\beta}_p
\]

\textbf{Fitted values:}
\[
\hat{Y}_i = \bhat{\beta}_0 + \bhat{\beta}_1 x_{i1} + \cdots + \bhat{\beta}_p x_{ip}
\]

Equivalently,
\[
Y_i = \beta_0 + \beta_1 x_{i1} + \cdots + \beta_p x_{ip} + \varepsilon_i
= \mathbf{x}_i^\top \bbeta + \varepsilon_i.
\]

Here,
\[
\bbeta = (\beta_0,\beta_1,\dots,\beta_p)^\top,
\qquad
\varepsilon_i \sim \mathcal{N}(0,\sigma^2) \text{ i.i.d.}
\]

The response variable $Y$ is a univariate random variable. The predictor vector is
\[
\mathbf{x} = (1, x_1, \dots, x_p)^\top,
\]
which is fixed and has dimension $p+1$.

\subsection{Model Properties}

\begin{enumerate}
    \item $\E(Y_i \mid X = \mathbf{x}_i) = \beta_0 + \beta_1 x_{i1} + \cdots + \beta_p x_{ip}$ (equivalent to $\E(\varepsilon_i) = 0$).

    \item On average, the response is linear in $\beta_0,\beta_1,\dots,\beta_p$ (linear in the parameters, not necessarily linear in the predictors).

    \item $\Var(Y_i \mid X = x_i) = \Var(Y_i \mid X_1 = x_{i1}, \ldots, X_p = x_{ip}) = \sigma^2$. This assumption means that the variability of the response around its conditional mean is constant for all values of the predictors (homoscedasticity).
\end{enumerate}

\subsection{Matrix Form of the Multiple Linear Regression Model}

The multiple linear regression model can be written as
\[
\mathbf{Y} = \mathbf{X}\bbeta + \beps,
\]
where
\[
\mathbf{X}
=
\begin{pmatrix}
1 & x_{11} & \cdots & x_{1p} \\
\vdots & \vdots & & \vdots \\
1 & x_{n1} & \cdots & x_{np}
\end{pmatrix},
\qquad
\bbeta
=
\begin{pmatrix}
\beta_0 \\
\vdots \\
\beta_p
\end{pmatrix},
\qquad
\beps
\sim N(\mathbf{0}_n, \sigma^2 \mathbf{I}_n).
\]

\subsection{Least Squares Estimator}

The parameters are estimated by the least squares estimator $\bhat{\beta}$ obtained by minimizing the residual sum of squares.

The fitted regression is
\[
\widehat{\E}(Y_i \mid X = \mathbf{x}_i)
= \hat{Y}_i
= \bhat{\beta}_0 + \bhat{\beta}_1 x_{i1} + \cdots + \bhat{\beta}_p x_{ip}
= \mathbf{x}_i^\top \bhat{\beta}.
\]

Combining all fitted values,
\[
\widehat{\E}(\mathbf{Y} \mid \mathbf{X})
= \hat{\mathbf{Y}}
= \mathbf{X}\bhat{\beta}
=
\begin{pmatrix}
\mathbf{x}_1^\top \bhat{\beta} \\
\vdots \\
\mathbf{x}_n^\top \bhat{\beta}
\end{pmatrix}.
\]

The least squares estimate $\bhat{\beta}$ of $\beta$ minimizes the residual sum of squares:
\[
\RSS(\beta)
= \sum_{i=1}^n \bigl(Y_i - \mathbf{x}_i^{\top}\beta\bigr)^2
= (\mathbf{Y} - \mathbf{X}\beta)^{\top}(\mathbf{Y} - \mathbf{X}\beta),
\]
where $Y_i$ is the observed value and $\mathbf{x}_i^{\top}\beta$ is the regression line value.

\subsection{Least Squares Procedure for Finding Estimators}

\begin{enumerate}
    \item Differentiate $\RSS(\beta)$ with respect to $\beta$ and set the derivative equal to $0$.
    \item Rearrange to solve for $\bhat{\beta}$.
    \item The minimizer is unique if $\mathbf{X}^{\top}\mathbf{X}$ is invertible.
\end{enumerate}

\subsection{Derivation of the Least Squares Estimator}

We expand the residual sum of squares:
\[
\begin{aligned}
\RSS(\beta)
&= (\bY - \mathbf{X}\beta)^{\top}(\bY - \mathbf{X}\beta) \\
&= (\bY^{\top} - \beta^{\top} \mathbf{X}^{\top})(\bY - \mathbf{X}\beta) \\
&= \bY^{\top}\bY - \bY^{\top}\mathbf{X}\beta - \beta^{\top}\mathbf{X}^{\top}\bY + \beta^{\top}\mathbf{X}^{\top}\mathbf{X}\beta.
\end{aligned}
\]

Since $\bY^{\top}\mathbf{X}\beta$ is a scalar,
\[
\bY^{\top}\mathbf{X}\beta = (\bY^{\top}\mathbf{X}\beta)^{\top} = \beta^{\top}\mathbf{X}^{\top}\bY,
\]
so
\[
\RSS(\beta)
= \bY^{\top}\bY - 2\beta^{\top}\mathbf{X}^{\top}\bY + \beta^{\top}\mathbf{X}^{\top}\mathbf{X}\beta.
\]

Differentiate with respect to $\beta$ (the first term is constant, the second is linear, the third is a quadratic form):
\[
\frac{\partial \RSS(\beta)}{\partial \beta}
= -2\mathbf{X}^{\top}\bY + 2\mathbf{X}^{\top}\mathbf{X}\beta.
\]

Setting the derivative equal to zero:
\[
-2\mathbf{X}^{\top}\bY + 2\mathbf{X}^{\top}\mathbf{X}\beta = 0,
\]
which simplifies to
\[
\mathbf{X}^{\top}\mathbf{X}\beta = \mathbf{X}^{\top}\bY.
\]

Assuming $\mathbf{X}^{\top}\mathbf{X}$ is invertible:
\[
\bhat{\beta} = (\mathbf{X}^{\top}\mathbf{X})^{-1}\mathbf{X}^{\top}\bY.
\]

The fitted values are therefore
\[
\hat{\bY} = \mathbf{X}\bhat{\beta} = \mathbf{X}(\mathbf{X}^{\top}\mathbf{X})^{-1}\mathbf{X}^{\top}\bY.
\]

\subsection{Mathematical Background}

Let $\mathbf{A}$ and $\mathbf{B}$ be matrices of sizes $m \times n$ and $n \times m$:
\[
(\mathbf{A}\mathbf{B})^{\top} = \mathbf{B}^{\top}\mathbf{A}^{\top}.
\]

For vectors $\mathbf{x}, \mathbf{y} \in \mathbb{R}^n$:
\[
\mathbf{x}^{\top}\mathbf{y} = \sum_{i=1}^n x_i y_i.
\]

\textbf{Linear function:} Let
\[
F(\beta) = \mathbf{a}^{\top}\beta = \sum_{j=0}^p a_j \beta_j,
\quad
\mathbf{a} = (a_0, a_1, \dots, a_p)^{\top}.
\]
Then
\[
\frac{\partial F(\beta)}{\partial \beta}
=
\begin{pmatrix}
\frac{\partial F}{\partial \beta_0} \\
\vdots \\
\frac{\partial F}{\partial \beta_p}
\end{pmatrix}
= \mathbf{a}.
\]

\textbf{Quadratic forms:} Consider
\[
Q(\beta) = \beta^{\top}\mathbf{Q}\beta,
\]
where $\mathbf{Q}$ is a $(p+1)\times(p+1)$ matrix. Then
\[
\frac{\partial Q(\beta)}{\partial \beta}
= (\mathbf{Q} + \mathbf{Q}^{\top})\beta.
\]

If $\mathbf{Q}$ is symmetric:
\[
\frac{\partial Q(\beta)}{\partial \beta} = 2\mathbf{Q}\beta.
\]

\begin{caution}
What assumptions did we make on the data for least squares estimation? Answer: $(\mathbf{X}^\top \mathbf{X})^{-1}$ must be invertible.
\end{caution}

\subsection{Example: Computing Parameter Estimates}

Given the data with sample size $n$ (to be determined from context):

\[
\sum y_i = 90,\quad
\sum x_{i,1} = 52,\quad
\sum x_{i,2} = 102,\quad
\sum x_{i,1}x_{i,2} = 536
\]

\[
\sum x_{i,1}^2 = 296,\quad
\sum x_{i,2}^2 = 1004,\quad
\sum x_{i,1}y_i = 482,\quad
\sum x_{i,2}y_i = 872
\]

Calculate the $\mathbf{X}^{\top}\mathbf{X}$ matrix:
\[
\mathbf{X}^\top \mathbf{X}
=
\begin{pmatrix}
n & \sum x_{i,1} & \sum x_{i,2} \\
\sum x_{i,1} & \sum x_{i,1}^2 & \sum x_{i,1}x_{i,2} \\
\sum x_{i,2} & \sum x_{i,2}x_{i,1} & \sum x_{i,2}^2
\end{pmatrix}
=
\begin{pmatrix}
n & 52 & 102 \\
52 & 296 & 536 \\
102 & 536 & 1004
\end{pmatrix}
\]

And the $\mathbf{X}^{\top}\mathbf{Y}$ vector:
\[
\mathbf{X}^\top \mathbf{Y}
=
\begin{pmatrix}
\sum y_i \\
\sum x_{i,1}y_i \\
\sum x_{i,2}y_i
\end{pmatrix}
=
\begin{pmatrix}
90 \\
482 \\
872
\end{pmatrix}
\]

Then:
\[
\bhat{\beta}
=
(\mathbf{X}^\top \mathbf{X})^{-1} \mathbf{X}^\top \mathbf{Y}
=
\begin{pmatrix}
0.9747634 & 0.2429022 & -0.2287066 \\
0.2429022 & 0.1620662 & -0.1111987 \\
-0.2287066 & -0.1111987 & 0.0835962
\end{pmatrix}
\begin{pmatrix}
90 \\
482 \\
872
\end{pmatrix}
\approx
\begin{pmatrix}
5.34 \\
2.96 \\
-1.30
\end{pmatrix}
\]

The fitted model is
\[
\hat{Y} = 5.34 + 2.96\,x_1 - 1.30\,x_2.
\]

\subsection{Important Distinction: Residuals vs.\ Errors}

\begin{notebox}
Let $\hat{e}_i = Y_i - \bhat{\beta}_0 - \bhat{\beta}_1 x_i$ be the residual (estimated from data) and $\varepsilon_i = Y_i - \beta_0 - \beta_1 x_i$ be the error (true but unobserved). These are different quantities: the first is an estimate, the second is a theoretical construct.
\end{notebox}

\subsection{Estimate of $\sigma^2$}

Residual sum of squares:
\[
\hat{e}_i
= Y_i - \hat{Y}_i
= Y_i - \mathbf{x}_i^\top \bhat{\beta},
\quad i = 1,\dots,n.
\]

\[
\RSS
= \sum_{i=1}^n \hat{e}_i^2
= \sum_{i=1}^n \left(Y_i - \mathbf{x}_i^\top \bhat{\beta}\right)^2
= (\mathbf{Y} - \mathbf{X}\bhat{\beta})^\top
   (\mathbf{Y} -\mathbf{X}\bhat{\beta}).
\]

Under the regression model assumption:
\[
\hat{\sigma}^2
= \frac{\RSS}{n - (p+1)}
\]
is used to estimate $\sigma^2$. The denominator is the number of observations $n$ minus the number of estimated parameters $(p+1)$. In some textbooks, $s^2$ is used to denote $\hat{\sigma}^2$.

\subsection{Interpretation of the Coefficients}

\begin{itemize}
    \item Fitted regression model:
    \[
    \widehat{\E}(Y \mid X=\mathbf{x})
    = \hat{Y}
    = \bhat{\beta}_0 + \bhat{\beta}_1 x_1 + \cdots + \bhat{\beta}_p x_p.
    \]

    \item $\bhat{\beta}_0$: estimated average of $Y \mid X=\mathbf{x}$ when all predictors are $0$.

    \item $\bhat{\beta}_1$: estimated average change in $Y \mid X=\mathbf{x}$ for a one-unit increase in $x_1$, holding all other predictors fixed.

    \item $\vdots$

    \item $\bhat{\beta}_p$: estimated average change in $Y \mid X=\mathbf{x}$ for a one-unit increase in $x_p$, holding all other predictors fixed.
\end{itemize}

\subsection{Categorical Predictors}

\begin{itemize}
    \item If a categorical predictor has $k$ categories (or levels) and the model has an intercept, then exactly $k - 1$ dummy variables are created.
    \item If a categorical predictor has $k$ categories (or levels) and the model does not have an intercept, then exactly $k$ dummy variables should be created.
\end{itemize}

\subsubsection*{Example: Design Matrix Dimensions with Categorical Predictor}

We have $n = 100$ observations with one response variable, one continuous predictor, and one categorical predictor with $k = 10$ levels, with an intercept included in the model. What is the dimension of the design matrix?

Since the categorical predictor has $10$ levels and the model includes an intercept, we create $k-1 = 9$ dummy variables. Therefore, the design matrix $\mathbf{X}$ has:
\[
\text{columns} = 1 \;(\text{intercept}) + 1 \;(\text{continuous predictor}) + 9 \;(\text{dummy variables}) = 11.
\]

Hence, the dimension of $\mathbf{X}$ is $\boxed{100 \times 11}$.

\subsubsection*{Example: Discrete vs.\ Continuous Predictor}

We have a discrete predictor that takes integer values from $1$ to $100$. Should it be treated as continuous or categorical?

This predictor should typically be treated as \textbf{continuous}, not categorical. Treating it as categorical would require $99$ dummy variables, which is inefficient. The values have a natural ordering and meaningful numerical scale. A continuous treatment preserves degrees of freedom and model interpretability.

\subsection{The Iris Data Example}

Measurements (in cm) of sepal length, sepal width, petal length, and petal width were collected for 50 flowers from each of three species of iris: \emph{setosa}, \emph{versicolor}, and \emph{virginica}. Our goal is to study the linear relationship between petal length (response) and the predictors: sepal length (continuous), sepal width (continuous), and species (categorical). We aim to predict petal length.

Consider predictor data of the form
\[
(A, 1), \ (A, -1), \ (C, 0.5), \ (B, 2).
\]

The fitted model using dummy variables is
\[
\widehat{\E}(Y_i \mid X = x_i)
= \hat{Y}_i
= \bhat{\beta}_0
+ \bhat{\beta}_1 I(x_i = B)
+ \bhat{\beta}_2 I(x_i = C).
\]

We create two dummy variables:
\begin{itemize}
    \item $I(x_i = B) = 1$ if the predictor category is $B$, and $0$ otherwise;
    \item $I(x_i = C) = 1$ if the predictor category is $C$, and $0$ otherwise.
\end{itemize}

Using an intercept and taking category $A$ as the reference level, the design matrix is

\[
\mathbf{X}
=
\begin{pmatrix}
1 & 0 & 0 \\
1 & 0 & 0 \\
1 & 0 & 1 \\
1 & 1 & 0
\end{pmatrix},
\qquad
\mathbf{Y}
=
\begin{pmatrix}
1 \\
-1 \\
0.5 \\
2
\end{pmatrix}.
\]

Assuming the fitted model using dummy variables with \emph{setosa} as the reference category is
\[
\widehat{\text{Petal}}_L
= -1.63
+ 2.17\,I(\text{versicolor})
+ 3.05\,I(\text{virginica})
+ 0.65\,\text{Sepal}_L
- 0.04\,\text{Sepal}_W.
\]

\subsubsection*{Example: Interpreting Categorical Coefficients}

\textbf{Question:} How do we interpret the coefficient for the species virginica? What are the dimensions of $\mathbf{X}$?

\textbf{Answer (a):} The coefficient $3.05$ represents the difference in the intercept between \emph{virginica} and the reference species \emph{setosa}, holding sepal length and sepal width fixed. For two flowers with the same sepal length and width, a \emph{virginica} flower is predicted to have a petal length that is $3.05$ cm longer than a \emph{setosa} flower on average.

Fitted equations for each species:

\textbf{Setosa:}
\[
\widehat{\text{Petal}}_L
= -1.63
+ 0.65\,\text{Sepal}_L
- 0.04\,\text{Sepal}_W.
\]

\textbf{Versicolor:}
\[
\widehat{\text{Petal}}_L
= (-1.63 + 2.17)
+ 0.65\,\text{Sepal}_L
- 0.04\,\text{Sepal}_W.
\]

\textbf{Virginica:}
\[
\widehat{\text{Petal}}_L
= (-1.63 + 3.05)
+ 0.65\,\text{Sepal}_L
- 0.04\,\text{Sepal}_W.
\]

\begin{notebox}
The intercepts for each category are:
\[
\bhat{\beta}_0 + \bhat{\beta}_2 = \text{Intercept for virginica},
\qquad
\bhat{\beta}_0 + \bhat{\beta}_1 = \text{Intercept for versicolor}.
\]
\end{notebox}

\textbf{Answer (b):} Dimensions of the Design Matrix. There are $n = 150$ observations in total. The predictors are: Intercept (1 column), Sepal length (1 column), Sepal width (1 column), Species with 3 levels $\Rightarrow$ $3 - 1 = 2$ dummy variables. Hence, the design matrix $\mathbf{X}$ has dimensions $150 \times 5$.

\clearpage
